\chapter{Przestrzeń Hilberta i notacja Diraca}

\begin{summarybox}
Dotychczas opisywaliśmy stany głównie za pomocą funkcji falowych \(\psi(x)\).
Formalizm Diraca pokazuje, że funkcja falowa jest tylko jedną reprezentacją
bardziej ogólnego obiektu: wektora stanu \(|\psi\rangle\) w przestrzeni Hilberta.
Ten rozdział wprowadza ket, bra, iloczyn skalarny, bazy ortonormalne, operatory,
projektory i związek między algebrą liniową a wcześniejszym opisem falowym.
\end{summarybox}

\section{Od funkcji falowej do wektora stanu}

W poprzednich rozdziałach stan cząstki opisywaliśmy najczęściej funkcją falową \(\psi(x)\). Taki opis jest bardzo konkretny: mówi, jak amplituda prawdopodobieństwa zależy od położenia. Nie jest jednak najbardziej ogólną postacią mechaniki kwantowej.

Bardziej podstawowym obiektem jest wektor stanu. Oznaczamy go symbolem Diraca:
\begin{equation}
    |\psi\rangle.
\end{equation}
Funkcja falowa jest reprezentacją tego wektora w bazie położeniowej:
\begin{equation}
    \psi(x)=\langle x|\psi\rangle.
\end{equation}
To równanie jest bardzo ważne. Mówi, że \(\psi(x)\) nie jest „rzeczą samą w sobie”, lecz współrzędną wektora stanu w konkretnej bazie.

Podobnie, reprezentacja pędowa ma postać
\begin{equation}
    \phi(p)=\langle p|\psi\rangle.
\end{equation}
Ten sam stan może więc mieć różne reprezentacje, tak jak ten sam wektor geometryczny może mieć różne współrzędne w różnych bazach.

\section{Ket \texorpdfstring{\(|\psi\rangle\)}{|psi>} i bra \texorpdfstring{\(\langle\psi|\)}{<psi|}}

W notacji Diraca wektor stanu zapisujemy jako ket:
\begin{equation}
    |\psi\rangle.
\end{equation}
Odpowiadający mu obiekt dualny zapisujemy jako bra:
\begin{equation}
    \langle\psi|.
\end{equation}
Jeżeli ket potraktujemy jak kolumnę, to bra jest sprzężonym hermitowsko wierszem. Dla wektora skończenie wymiarowego
\begin{equation}
    |\psi\rangle=\begin{pmatrix}c_1\\ c_2\\ \vdots\\ c_n\end{pmatrix},
\end{equation}
bra ma postać
\begin{equation}
    \langle\psi|=\begin{pmatrix}c_1^* & c_2^* & \cdots & c_n^*\end{pmatrix}.
\end{equation}

Zapis bra-ket jest wygodny, ponieważ od razu rozróżnia wektor, jego dual oraz iloczyn skalarny. Wyrażenie
\begin{equation}
    \langle\phi|\psi\rangle
\end{equation}
jest liczbą zespoloną, natomiast
\begin{equation}
    |\phi\rangle\langle\psi|
\end{equation}
jest operatorem.

\section{Iloczyn skalarny}

Iloczyn skalarny dwóch stanów zapisujemy jako
\begin{equation}
    \langle\phi|\psi\rangle.
\end{equation}
W przestrzeni funkcji falowych odpowiada mu całka
\begin{equation}
    \langle\phi|\psi\rangle=\int \phi^*(x)\psi(x)\,dx.
\end{equation}
Norma stanu wynosi
\begin{equation}
    \langle\psi|\psi\rangle.
\end{equation}
Stan znormalizowany spełnia
\begin{equation}
    \langle\psi|\psi\rangle=1.
\end{equation}

Iloczyn skalarny ma interpretację amplitudy przejścia. Liczba \(\langle\phi|\psi\rangle\) mówi, jaka jest amplituda znalezienia stanu \(|\psi\rangle\) w kierunku stanu \(|\phi\rangle\). Prawdopodobieństwo wynosi
\begin{equation}
    P_{\phi}=|\langle\phi|\psi\rangle|^2.
\end{equation}

\begin{remarkbox}
Kolejność w iloczynie skalarnym ma znaczenie. Zwykle \(\langle\phi|\psi\rangle\) jest sprzężeniem zespolonym \(\langle\psi|\phi\rangle\).
\end{remarkbox}

\section{Baza ortonormalna}

Baza ortonormalna to zbiór stanów \(|n\rangle\), które spełniają warunek
\begin{equation}
    \langle m|n\rangle=\delta_{mn}.
\end{equation}
Każdy stan można wtedy rozwinąć w tej bazie:
\begin{equation}
    |\psi\rangle=\sum_n c_n|n\rangle.
\end{equation}
Współczynniki rozwinięcia otrzymujemy przez rzutowanie:
\begin{equation}
    c_n=\langle n|\psi\rangle.
\end{equation}

Warunek zupełności bazy zapisujemy jako
\begin{equation}
    \sum_n |n\rangle\langle n|=\hat{I}.
\end{equation}
To jest wersja operatorowa stwierdzenia, że baza obejmuje całą przestrzeń stanów.

Dla zmiennej ciągłej, na przykład położenia, zamiast sumy mamy całkę:
\begin{equation}
    \int |x\rangle\langle x|\,dx=\hat{I}.
\end{equation}
Wtedy funkcja falowa \(\psi(x)=\langle x|\psi\rangle\) jest współczynnikiem rozwinięcia stanu w bazie położeniowej.

\section{Rozkład stanu w bazie}

Jeżeli stan jest rozwinięty jako
\begin{equation}
    |\psi\rangle=\sum_n c_n|n\rangle,
\end{equation}
to normalizacja oznacza
\begin{equation}
    \sum_n |c_n|^2=1.
\end{equation}
Liczba \(|c_n|^2\) jest prawdopodobieństwem otrzymania wyniku związanego ze stanem bazowym \(|n\rangle\), jeśli baza jest bazą stanów własnych mierzonej obserwabli.

Dla układu dwupoziomowego, na przykład kubitu, stan ma postać
\begin{equation}
    |\psi\rangle=\alpha|0\rangle+\beta|1\rangle,
\end{equation}
przy warunku
\begin{equation}
    |\alpha|^2+|\beta|^2=1.
\end{equation}
Ten przykład stanie się podstawą następnego rozdziału.

\section{Operatory jako macierze}

Operator \(\hat{A}\) działa na ket i zwraca nowy ket:
\begin{equation}
    \hat{A}|\psi\rangle=|\phi\rangle.
\end{equation}
W danej bazie operator ma reprezentację macierzową. Elementy macierzy są dane wzorem
\begin{equation}
    A_{mn}=\langle m|\hat{A}|n\rangle.
\end{equation}
Jeżeli
\begin{equation}
    |\psi\rangle=\sum_n c_n|n\rangle,
\end{equation}
to współczynniki stanu \(\hat{A}|\psi\rangle\) otrzymujemy przez zwykłe mnożenie macierzy przez wektor.

Wartość oczekiwana obserwabli \(A\) ma w notacji Diraca postać
\begin{equation}
    \langle\hat{A}\rangle=\langle\psi|\hat{A}|\psi\rangle.
\end{equation}
To jest ten sam wzór, który wcześniej zapisywaliśmy jako całkę
\begin{equation}
    \int \psi^*(x)\hat{A}\psi(x)\,dx.
\end{equation}

\section{Projektory i pomiar}

Projektor na stan \(|a\rangle\) definiujemy jako
\begin{equation}
    \hat{P}_a=|a\rangle\langle a|.
\end{equation}
Jeżeli stan \(|\psi\rangle\) jest znormalizowany, to prawdopodobieństwo otrzymania wyniku odpowiadającego \(|a\rangle\) wynosi
\begin{equation}
    p(a)=\langle\psi|\hat{P}_a|\psi\rangle=|\langle a|\psi\rangle|^2.
\end{equation}
Po idealnym pomiarze, który dał wynik \(a\), stan zostaje związany z odpowiednim podprzestrzennym projektorem. Dla nieosobliwego przypadku jednowymiarowego oznacza to przejście do stanu \(|a\rangle\), z dokładnością do fazy.

Projektory są bardzo wygodne, bo łączą algebrę liniową z regułą Borna. Pomiar nie jest wtedy osobnym, tajemniczym dodatkiem do formalizmu, lecz operacją rzutowania na odpowiednie podprzestrzenie.

\section{Związek z wcześniejszymi funkcjami falowymi}

Funkcja falowa w położeniu jest po prostu reprezentacją ketu w bazie \(|x\rangle\):
\begin{equation}
    \psi(x)=\langle x|\psi\rangle.
\end{equation}
Normalizacja
\begin{equation}
    \int |\psi(x)|^2\,dx=1
\end{equation}
jest tym samym co
\begin{equation}
    \langle\psi|\psi\rangle=1.
\end{equation}
Wartość oczekiwana położenia
\begin{equation}
    \langle x\rangle=\int x|\psi(x)|^2\,dx
\end{equation}
jest reprezentacją wzoru Diraca
\begin{equation}
    \langle\hat{x}\rangle=\langle\psi|\hat{x}|\psi\rangle.
\end{equation}

Notacja Diraca nie zastępuje funkcji falowych, lecz pokazuje, że są one jednym z wielu możliwych zapisów tego samego stanu.

\section{Mathematica jako algebra liniowa mechaniki kwantowej}

W skończenie wymiarowych przykładach Mathematica może traktować stany jako wektory, a operatory jako macierze. Przykładowo:
\begin{verbatim}
psi = {alpha, beta};
A = {{a11, a12}, {a21, a22}};
Conjugate[psi].psi
Conjugate[psi].A.psi
\end{verbatim}
Dla wektorów zespolonych lepiej używać sprzężenia hermitowskiego:
\begin{verbatim}
ConjugateTranspose[{psi}].A.{psi}
\end{verbatim}
W praktyce trzeba pilnować kształtu obiektów: czy wektor jest listą, kolumną, czy macierzą jednowierszową. Dla celów dydaktycznych często najczytelniej jest pisać małe funkcje pomocnicze:
\begin{verbatim}
bra[v_] := Conjugate[v]
expect[A_, v_] := bra[v].A.v
\end{verbatim}

\begin{summarybox}
Notacja Diraca porządkuje mechanikę kwantową jako algebrę liniową. Ket jest stanem, bra jest obiektem dualnym, iloczyn bra-ket jest amplitudą, operator działa na stan, a funkcja falowa jest reprezentacją ketu w wybranej bazie.
\end{summarybox}
